% Direct-authority diplomatic French for printed p.138.
\SourcePage{138}

\LayerTarget{tohoku:def:subcategory:p2}
\noindent plus des conditions précédentes, les groupes
$\mathrm{Hom}_{C'}(A,B)$ sont des sous-groupes de
$\mathrm{Hom}_C(A,B)$.

\LayerTarget{tohoku:def:complete-subcategory}
Supposons que $C$ soit une catégorie abélienne, $C'$ est appelé une
\textit{sous-catégorie complète} si (i) \textit{pour $A,B\in C'$, on a
$\mathrm{Hom}_{C'}(A,B)=\mathrm{Hom}_C(A,B)$.} (ii) \textit{Si dans une
suite exacte $A\longrightarrow B\longrightarrow C\longrightarrow
D\longrightarrow E$ les 4 termes extrêmes sont dans $C'$, il en est de même
du terme médian $D$.} D'après (i), $C'$ est complètement déterminé par la
classe de ses objets. (ii) équivaut à dire que pour tout morphisme
$P\longrightarrow Q$ avec $P,Q\in C'$ le noyau et le conoyau sont dans $C'$,
et que pour toute suite exacte
$0\longrightarrow R'\longrightarrow R\longrightarrow R''\longrightarrow 0$
avec $R',R''$ dans $C'$, $R$ est dans $C'$. On vérifie aussitôt, qu'alors
$C'$ est elle-même une classe abéliennne, et que pour un morphisme
$u:A\longrightarrow B$ dans $C'$, noyau et conoyau (donc aussi image et
coimage) sont identiques au noyau et conoyau (resp. \ldots) envisagé dans
$C$.

\LayerTarget{tohoku:def:thick-subcategory}
La sous-catégorie $C'$ de $C$ est dite \textit{épaisse}, si elle satisfait à
la condition (i) ci-dessus, et au renforcement suivant de la condition (ii) :
(iii) \textit{Si dans une suite exacte
$A\longrightarrow B\longrightarrow C$ les termes extrêmes $A,C$ sont dans
$C'$, il en est de même de $B$.} Si $C$ est la catégorie abélienne des
groupes abéliens, on trouve la notion de ``classe de groupes abéliens'' de
[17]. On voit comme dans [17] que (iii) équivaut à la conjonction des trois
conditions suivantes : Tout objet nul est dans $C'$, tout objet isomorphe à
un sous-truc ou un truc-quotient d'un objet $C'$ est dans $C'$, toute
extension de deux objets de $C'$ est dans $C'$.

\LayerTarget{tohoku:def:quotient-category}
Soit $C$ une catégorie abélienne, $C'$ une sous-catégorie épaisse, nous allons
définir une nouvelle catégorie abélienne, notée $C/C'$ et appelée
\textit{catégorie quotient de $C$ par $C'$}. Les objets de $C/C'$ sont par
définition les objets de $C$. Nous allons définir les morphismes dans $C/C'$
de $A$ dans $B$, appelés ``morphismes mod. $C'$ de $A$ dans $B$''.

\LayerTarget{tohoku:def:quasi-equal}
Disons qu'un sous-truc $A'$ de $A$ est égal mod. $C'$ à $A$, ou quasi-égal à
$A$, si $A/A'\in C'$ ; alors tout sous-truc de $A$ contenant $A'$ est encore
quasi-égal à $A$, de plus le inf de deux sous-trucs de $A$ quasi-égaux à $A$
est encore quasi-égal à $A$. Dualement, on introduit la notion de quotient de
$B$ qussi-égal à $B$ : un tel quotient $B/N$ est quasi-égal à $B$ si
$N\in C'$.

\LayerTarget{tohoku:def:quotient-morphism}
Un morphisme mod. $C'$ de $A$ dans $B$ est alors défini par un morphisme $f'$
d'un sous-truc $A'$ de $A$ quasi-égal à $A$, dans un quotient $B'$ de $B$
quasi-égal à $B$, étant entendu qu'un morphisme
$f'':A''\longrightarrow B''$ (satisfaisant aux mêmes conditions), définit le
même morphisme mod. $C'$, si et seulement si on peut trouver
$A'''\leq\mathrm{Inf}(A',A'')$, $B'''\leq\mathrm{Inf}(B',B'')$, $A'''$ quasi
égal à $A$, $B'''$ quasi égal à $B$, tels que les morphismes
$A'''\longrightarrow B'''$ définis par $f',f''$ soient identiques. Cette
dernière relation entre $f'$ et $f''$ est bien une relation d'équivalence, la
définition précédente des morphismes mod. $C'$ est donc cohérente.

\LayerTarget{tohoku:def:quotient-hom}
Supposons que pour tout $A\in C$, les sous-trucs de $A$ forment un ensemble
(ce qui est vrai pour toutes les catégories connues), alors on peut
considérer l'ensemble des morphismes mod. $C'$ de $A$ dans $B$, noté
$\mathrm{Hom}_{C/C'}(A,B)$. $\mathrm{Hom}_{C/C'}(A,B)$ apparaît comme limite
inductive des groupes abéliens $\mathrm{Hom}_C(A',B')$ où $A'$ ($B'$)
parcourt les sous-trucs de $A$ (quotients de $B$) quasi-égaux à $A$ ($B$), et
peut être considéré par suite comme un groupe abélien.

\LayerTarget{tohoku:def:quotient-composition:p1}
On définit de même un accouplement

\[
 \mathrm{Hom}_{C/C'}(A,B)\times\mathrm{Hom}_{C/C'}(B,C)
 \longrightarrow\mathrm{Hom}_{C/C'}(A,C)
\]

\noindent de la façon suivante. Soit $u\in\mathrm{Hom}_C(A',B')$ un
représentant d'une $u'\in\mathrm{Hom}_{C/C'}(A,B)$, et
$v\in\mathrm{Hom}_C(B'',C')$ un représentant d'un
$v'\in\mathrm{Hom}_{C/C'}(B,C)$. Soit $Q$ l'image du morphisme canonique
$B''\longrightarrow B'$ du

\TohokuContinuation{139}
